%!TEX root = ../main.tex
% Sources: ~/WorkEnv/ai/tests/supervised/{sine_*,lnn_*}/*_metrics.csv,
%          ~/WorkEnv/ai/tests/rl/lunar_*/training_log.csv,
%          corpus statistics from build_manifest(~/WorkEnv/core/datasets),
%          figures from ~/WorkEnv/ai/plots/*/{sine_all,lnn_all,lunar_all} -> images/ch08/.
\chapter{Results}
\label{ch:results}

\section{Supervised Learning}
\label{sec:res-sl}

\subsection{Validation on a Synthetic Task}
\label{sec:res-sl-synthetic}

The first result concerns the machinery rather than the problem. Before any
network was trained on the corpus, the seven vector twins of
Section~\ref{sec:networks-implementation} --- the six recurrent networks and the
stateless control, each with its cells and wiring but without the encoder ---
were trained on the synthetic task of Section~\ref{sec:exp-protocol}: predicting
the next sample of a sine wave whose phase is drawn uniformly in $[0, 2\pi)$,
whose frequency is drawn in $[0.8, 1.2]$, and to which Gaussian noise of
standard deviation $0.05$ is added. The dataset holds \num{1000} waves of
\num{50} steps, split 80/20 between training and validation. Every session lasts
\num{30} epochs, with an initial learning rate of $10^{-3}$ decaying linearly to
15\% of that value and batches of \num{32} sequences, that is \num{25}
optimisation steps per epoch. All sessions share the same seed, for the weights
and for the data.

\paragraph{The references.}
The task has two values known in advance, and that is what makes it a check
rather than an experiment. The first is the error of the zero predictor: since
the wave has zero mean and unit amplitude, its mean energy is $\tfrac12$, and
the \met{val_baseline_zero_loss} column indeed reads $0.498$--$0.505$ in every
session. The second is the noise floor: the sample to be predicted carries
independent noise of variance $0.05^2 = 0.0025$, which no model can anticipate.
The largest attainable coefficient of determination is therefore about
$1 - 0.0025/0.5 \approx 0.995$, and every result is to be read as a distance
from that floor, not from zero.

\paragraph{Outcome.}
Table~\ref{tab:res-sine} summarises the seven sessions at the last epoch, which
for all of them is also the best in validation: the recurrent networks are still
improving after thirty epochs, and the control has been flat since the tenth.

\begin{table}[htbp]
  \centering
  \small
  \caption{Synthetic task: next-step prediction of a noisy sine wave.
  Validation values at epoch 30 (the best one for every session). The noise
  floor is $0.0025$, the zero predictor about $0.50$. \emph{Epoch} is the first
  epoch at which \texttt{val\_loss} falls below $0.05$.}
  \label{tab:res-sine}
  \begin{tabular}{@{}lrrrrr@{}}
    \toprule
    \textbf{Network} & \met{val_loss} & $R^2$ & \textbf{Amplitude} & \textbf{Epoch} & \textbf{s/epoch} \\
    \midrule
    \met{LNN_CfC_NCP}  & 0.0061 & 0.988 & 1.000 & 1  & 1.50 \\
    \met{LNN_CfC}      & 0.0072 & 0.986 & 0.992 & 1  & 1.34 \\
    \met{LNN_LTC_NCP}  & 0.0125 & 0.975 & 0.984 & 5  & 2.37 \\
    \met{LNN_LRCU_NCP} & 0.0180 & 0.964 & 0.977 & 6  & 4.17 \\
    \met{LNN_LTC}      & 0.0184 & 0.963 & 0.978 & 5  & 1.95 \\
    \met{MLP} (control) & 0.0363 & 0.928 & 0.960 & 1 & 0.12 \\
    \met{LNN_LRCU}     & 0.1398 & 0.722 & 0.845 & -- & 3.22 \\
    \bottomrule
  \end{tabular}
\end{table}

Six networks out of seven beat the control, and the two CfC variants reach
$0.006$--$0.007$, less than three times the noise floor, with an amplitude ratio
of $1.00$ and $0.99$: the prediction has the same spread as the target and shows
no regression towards the mean. The recurrent path --- windows, state
propagation, truncated backpropagation --- therefore works end to end for all
three cells, which is what the task was meant to establish.

\paragraph{The stateless control stops where it should.}
The value at which the \met{MLP} stops is not a defect of the training but a
consequence of the task, and being able to predict it is the most useful check
the task provides. A network without memory sees only the current sample
$x_t = \sin\theta$ and cannot tell whether the wave is rising or falling: the
sign of $\cos\theta$ is not available to it. The best predictor it can realise
is then $x_t\cos\omega$, with $\omega \approx 4\pi/50$ the phase advance per
step, whose expected error is $\tfrac12\sin^2\omega \approx 0.031$; adding the
noise on the target and the noise propagated from the input gives about
$0.036$. A numerical check --- the conditional mean of the next sample given the
current one, estimated on \num{20000} freshly generated waves --- gives
$0.0359$. The \met{MLP} curve flattens at $0.0363$ from the tenth epoch onwards
(Figure~\ref{fig:res-sine-losses}). The control, in other words, measures
exactly what memory is worth on this task, and everything the recurrent
networks gain below that value is information extracted from the past of the
sequence.

\paragraph{Speed of convergence.}
The cells differ more sharply in \emph{when} than in \emph{how much}. The two
CfC networks and the control are below $0.05$ after the first epoch. LTC and
LRCU instead spend their first two or three epochs at $\approx 0.5$, that is on
the value of the zero predictor, and leave it only between the third and the
sixth. The \met{LNN_LRCU} without NCP wiring does not leave it entirely: after
thirty epochs it stands at $0.14$ ($R^2 = 0.72$), with a curve that on a
logarithmic scale is still descending steadily. It is to be read as a network
that has not converged within the step budget it was given, not as a cell
incapable of the task: its NCP variant, which differs only in the wiring and in
the two skip connections, reaches $0.018$.

\paragraph{No overfitting, and nothing for the calibration to do.}
The ratio between validation and training loss settles around one from the
tenth epoch onwards for every network (Figure~\ref{fig:res-sine-losses},
bottom): with a thousand waves drawn from the same distribution there is nothing
to memorise. The output calibration (Section~\ref{sec:output-calibration}) has
nothing to correct either: \met{val_cal_loss} agrees with \met{val_loss} to the
third significant figure in every session, as it must when the amplitude is
already one. The observation matters by contrast with the corpus, where the
calibration does intervene.

\paragraph{State norms.}
The only anomaly the task hints at concerns LRCU. The mean norm of the first
cell's state grows during training from $2.4$ to $7.9$ for the \met{LNN_LRCU}
and from $2.4$ to $6.5$ for its NCP variant, whereas for CfC it stays between
$1.8$ and $3.0$ and for LTC between $2.1$ and $2.5$. Over fifty steps the growth
is harmless; over episodes of a thousand steps it will not be
(Section~\ref{sec:res-rl-lunar}).

\begin{figure}[htbp]
  \centering
  \includegraphics[width=\textwidth]{ch08/sine_losses.png}
  \caption{Synthetic task: training and validation loss on a linear and a
  logarithmic scale (top), and the gap and ratio between validation and
  training (bottom), for the seven vector networks.}
  \label{fig:res-sine-losses}
\end{figure}

\paragraph{What the task does not verify.}
The time step is constant ($\Delta t = 1$), so the task exercises neither the
time argument of the CfC cells nor the interval appended to the observation for
LTC and LRCU (Section~\ref{sec:cell-time}). The task also has a single output
and knows no velocity caps: the cap-scaled loss falls back here to the plain
squared error. What it establishes is that each cell, wired as it will be on the
corpus, learns a known dynamics; everything else is left to the later levels.


\subsection{Training on the Recorded Corpus}
\label{sec:res-sl-corpus}

\paragraph{Setting.}
The seven spatial networks were trained on the recorded corpus with the
run-stratified split of Section~\ref{sec:sl-split}: \num{332} episodes
(\num{77240} frames) for training, \num{74} (\num{16009}) for validation and
\num{70} (\num{17747}) for test, stratified on robot model, world and voxel
size. The corpus contains two robot models --- an unmanned ground vehicle (UGV)
and an unmanned aerial vehicle (UAV) --- and two worlds, and the three
partitions hold them in similar proportions. Every session lasts \num{20}
epochs on windows of \num{16} steps with a stride of \num{4}, batches of
\num{8} sequences, and a learning rate going from $10^{-4}$ to
$1.5\cdot10^{-5}$ with linear decay; an epoch counts \num{1340} measured steps
and lasts between \num{557} and \num{601} seconds, about three and a quarter
hours per session. The remaining settings, read back from the checkpoints'
metadata, are also common to the seven sessions: weights seed and split seed
both $0$, \num{64} archives open at once for the composition of the batch
(Section~\ref{sec:sl-sequences}) with a shuffling reservoir of \num{64}
batches, gradient clipping at a norm of $5$, no position features, and the
auxiliary geometry head enabled with a weight of $0.3$
(Section~\ref{sec:sl-geometry}). Training windows are augmented with the
left--right reflection of Section~\ref{sec:sl-augmentation}.

This subsection describes what the metrics say about the imitation as such,
taking as reference the network that is best in validation,
\met{SparseConv_CfC_NCP}; the comparison between architectures is deferred to
Section~\ref{sec:res-sl-comparison}. All values are read on the epoch selected
as prescribed in Section~\ref{sec:exp-protocol}, that is the one with the
lowest \met{val_loss}. The test partition, read once on the selected
checkpoints, is reported at the end of Section~\ref{sec:res-sl-comparison}.

\paragraph{How the energy of the command is distributed.}
Before reading any error it pays to look at the target. Of the six components
of the command, two --- the roll and pitch rates --- are identically zero across
the whole corpus, and their $R^2$ is undefined. Of the remaining four, the
forward velocity alone carries about 86\% of the energy of the command over the
476 episodes (\num{110996} frames) the sessions were run on, and 84\% and 83\%
in the training and validation partitions respectively (measured as the error
of the zero predictor, the \met{baseline_dim_i} columns); the vertical and yaw
rates share the rest almost equally, about 7\% each, and the lateral velocity
carries less than $0.1\%$. It is this concentration that makes the per-component reading of
Section~\ref{sec:metrics-r2} necessary: a mean loss can be determined almost
entirely by the forward component alone.

The second property of the target that conditions every reading is the share of
frames whose command is identically zero on every component. Over the corpus
it is $12.4\%$, and it is essentially the same for the two platforms
($12.6\%$ for the UGV, $12.2\%$ for the UAV). They are legitimate ``hold still''
targets rather than missing labels, which is why the error on them is reported
separately (Section~\ref{sec:metrics-still}).

\paragraph{The loss against the references.}
The \met{SparseConv_CfC_NCP} reaches its lowest \met{val_loss} at epoch \num{11}
with $4.63\cdot10^{-3}$, that is $3.7\%$ of the zero predictor ($0.126$) and
$35\%$ of the predictor that echoes the observed velocity
($1.31\cdot10^{-2}$). The second comparison is the one that matters: echoing
the velocity is a strong predictor because the command changes slowly, and
beating it by a factor of three says that the network is using the observation
and not merely the velocity the observation already contains.

Read per component, the same network has $R^2 = 0.975$ on the forward velocity,
$0.766$ on the vertical one and $0.982$ on yaw
(Figure~\ref{fig:res-corpus-r2}). The skill against the echo, however, tells a
different story for yaw: it is $0.75$ forward, $0.55$ vertically and
\emph{negative}, $-0.17$, on yaw. The yaw rate in the corpus is so persistent
that echoing it yields a lower error ($3.4\cdot10^{-4}$) than the network's
($4.0\cdot10^{-4}$): the high $R^2$ on that component is to the credit of the
signal's persistence, not of the model. None of the seven networks beats the
echo on yaw (Figure~\ref{fig:res-corpus-skill}).

The lateral component is the opposite case. With less than $0.1\%$ of the
energy, its raw $R^2$ is negative ($-1.6$ for the \met{SparseConv_CfC_NCP},
between $-0.08$ and $-1.6$ for the others): the network produces a lateral
command that is small in absolute terms but large relative to a target that is
almost always zero. The calibration brings the calibrated $R^2$ back to $0.00$
and the amplitude ratio to $0.03$, that is, it effectively switches the axis
off. This is the correct behaviour given the corpus, and it is to be read as an
absence of data on that component, not as learning.

\begin{figure}[htbp]
  \centering
  \includegraphics[width=\textwidth]{ch08/corpus_r2_per_dim.png}
  \caption{Corpus: coefficient of determination per component of the command,
  for the seven spatial networks. Components 3 and 4 (roll and pitch rate) are
  zero in the corpus and have no defined $R^2$.}
  \label{fig:res-corpus-r2}
\end{figure}

\begin{figure}[htbp]
  \centering
  \includegraphics[width=\textwidth]{ch08/corpus_skill_vs_echo.png}
  \caption{Corpus: skill against the predictor that echoes the observed
  velocity, per component. Positive values mean the model beats the echo.}
  \label{fig:res-corpus-skill}
\end{figure}

\paragraph{Amplitude and restart.}
Over all frames the regression towards the mean is limited: the amplitude ratio
of the \met{SparseConv_CfC_NCP} is $0.97$ forward, $0.81$ vertically and $0.97$
on yaw, and the calibration raises the vertical one to $0.88$
(Figure~\ref{fig:res-corpus-amplitude}). The picture changes completely on the
restart frames, where the robot is at rest and the demonstration starts moving.
There the forward restart ratio is $0.02$: the network commands $2\%$ of the
velocity the pilot commanded ($7\cdot10^{-4}$ against
$3.5\cdot10^{-2}$\,\si{\metre\per\second}). Across the seven networks the value
lies between $0.01$ and $0.09$, and the calibration does not move it, because it
corrects the overall spread and not a rare event. On the vertical component the
ratio is higher, between $0.17$ and $0.25$, but it remains far from one.

This is exactly the failure anticipated in
Section~\ref{sec:metrics-amplitude}: at the moment of departure the
demonstration is ambiguous --- the same state of rest sometimes precedes a
departure and sometimes a further wait --- and the least-squares optimum is to
command almost nothing. No architecture solves it by imitation, and it is one of
the reasons why the second training stage exists.

\begin{figure}[htbp]
  \centering
  \includegraphics[width=\textwidth]{ch08/corpus_amplitude.png}
  \caption{Corpus: amplitude ratio and restart ratio per component, raw and
  calibrated.}
  \label{fig:res-corpus-amplitude}
\end{figure}

\paragraph{What the selected epoch gives up.}
The protocol of Section~\ref{sec:exp-protocol} selects the epoch of lowest
validation loss, but asks that an epoch improving the loss while degrading the
amplitude or restart metrics be flagged rather than selected in silence. On the
amplitude there is nothing to flag: for every network the selected epoch is
also the one with the largest calibrated amplitude ratio on the forward
component, or within $0.003$ of it. On the restart ratio every network has to
be flagged. The ratio is highest in the first epochs and falls as the loss
falls: from $0.39$ at epoch~1 to $0.06$ at the selected epoch for the
\met{SparseConv_LRCU}, from $0.27$ to $0.07$ for the \met{SparseConv_LTC},
from $0.25$ to $0.09$ for the \met{SparseConv_LRCU_NCP} and from $0.11$ to
$0.01$ for the \met{SparseConv_CfC_NCP} (calibrated values, which are what the
robot executes). Training, in other words, buys part of its loss by commanding
less at departure, which is the behaviour a squared error rewards. The epochs
that command more at a restart have a higher loss --- between $1.4$ and $3.9$
times the selected one --- and choosing them would trade one defect for
another. One exception deserves to be named: for the
\met{SparseConv_LTC_NCP}, epoch~15 has a restart ratio of $0.13$ against
$0.08$ at the selected epoch~20, for a validation loss only 3\% higher. The
selection was nevertheless kept on the loss, as the protocol prescribes; how
much the restart weighs on the fleet is the subject of
Section~\ref{sec:res-cross}.

\paragraph{Still and moving.}
One would expect the frames with a zero command to be the easiest. The opposite
happens: for the \met{SparseConv_CfC_NCP} the per-frame error on still frames
is $5.8\cdot10^{-3}$, against $3.0\cdot10^{-3}$ on moving ones, and for every
network the error on still frames exceeds the error on moving ones, by a factor
between $1.3$ and $2.2$ (Figure~\ref{fig:res-corpus-still}). When the pilot
holds still, the network keeps commanding a small residual velocity. Read
together with the restart ratio, the picture is consistent: around rest the
network errs in both directions, commanding too little when the pilot departs
and a little when the pilot holds still --- it interpolates between the two
cases instead of telling them apart.

\begin{figure}[htbp]
  \centering
  \includegraphics[width=\textwidth]{ch08/corpus_still_vs_moving.png}
  \caption{Corpus: error on frames with a zero command and on frames in
  motion.}
  \label{fig:res-corpus-still}
\end{figure}

\paragraph{Per robot.}
The per-robot breakdown (Section~\ref{sec:metrics-per-robot}) shows neither
platform served worse than the other on the components they share: on the
forward velocity the \met{SparseConv_CfC_NCP} has $R^2 = 0.971$ on the UGV and
$0.978$ on the UAV, on yaw $0.986$ and $0.976$. Only the UAV commands the
vertical component, and the figures on it coincide in practice with the
aggregate ones; the lateral component, which the UGV cannot command, is
practically absent from the UAV's demonstrations as well.

\paragraph{Overfitting.}
On the corpus, unlike on the sine wave, the model memorises. For the
\met{SparseConv_CfC_NCP} the training loss keeps falling until the last epoch,
from $2.9\cdot10^{-3}$ at epoch 11 to $1.6\cdot10^{-3}$ at epoch 20, while the
validation loss stops improving at epoch 11 and then oscillates between $4.8$
and $6.1\cdot10^{-3}$. The ratio between the two rises to about three by the
twentieth epoch (Figure~\ref{fig:res-corpus-losses}). The behaviour is
consistent with a corpus of \num{332} training episodes and with networks sized
for the problem rather than for the corpus
(Section~\ref{sec:networks-implementation}): the excess capacity is spent
memorising the episodes seen. It is also why selecting the checkpoint on the
validation split is not a detail.

\begin{figure}[htbp]
  \centering
  \includegraphics[width=\textwidth]{ch08/corpus_losses.png}
  \caption{Corpus: training and validation loss (top) and their gap and ratio
  (bottom), for the seven spatial networks.}
  \label{fig:res-corpus-losses}
\end{figure}


\subsection{Architecture Comparison}
\label{sec:res-sl-comparison}

Table~\ref{tab:res-corpus} compares the seven networks at the epoch of lowest
\met{val_loss}, under the conditions of Section~\ref{sec:exp-protocol}: same
corpus, same split, same seed, same schedule, same state widths.

\begin{table}[htbp]
  \centering
  \small
  \caption{Corpus: comparison of the architectures at the epoch of lowest
  validation loss. Losses in units of $10^{-3}$; $R^2$, skill and restart
  ratio on the forward component ($v_x$) unless stated otherwise, on raw
  values. Zero predictor $126.4$, velocity echo $13.1$.}
  \label{tab:res-corpus}
  \setlength{\tabcolsep}{4pt}
  \begin{tabular}{@{}lrrrrrrrr@{}}
    \toprule
    \textbf{Network} & \textbf{Ep.} & \met{val_loss} & \textbf{cal.} & $R^2_{v_x}$ & $R^2_{v_z}$ & $R^2_{\omega_z}$ & \textbf{Skill} & \textbf{Restart} \\
    \midrule
    \met{CfC_NCP}  & 11 &  4.63 &  4.49 & 0.975 & 0.766 & 0.982 & 0.75 & 0.02 \\
    \met{MLP} (ctrl.) & 11 &  5.09 &  4.92 & 0.969 & 0.766 & 0.982 & 0.69 & 0.03 \\
    \met{CfC}      & 11 &  5.52 &  5.35 & 0.965 & 0.752 & 0.981 & 0.65 & 0.01 \\
    \met{LTC_NCP}  & 20 &  8.67 &  8.58 & 0.948 & 0.582 & 0.969 & 0.49 & 0.09 \\
    \met{LRCU_NCP} & 20 &  9.98 &  9.88 & 0.939 & 0.562 & 0.959 & 0.40 & 0.08 \\
    \met{LTC}      & 20 & 11.78 & 11.74 & 0.936 & 0.534 & 0.926 & 0.37 & 0.08 \\
    \met{LRCU}     & 20 & 13.03 & 12.87 & 0.933 & 0.453 & 0.919 & 0.33 & 0.06 \\
    \bottomrule
  \end{tabular}
\end{table}

\paragraph{Two groups.}
The networks split sharply into two groups. The two CfC networks and the control
reach their minimum at epoch 11 with losses between $4.6$ and
$5.5\cdot10^{-3}$, and then overfit. The four LTC and LRCU networks are instead
still descending at the last epoch, with losses $1.9$ to $2.8$ times higher,
and the \met{SparseConv_LRCU} barely matches the velocity echo ($13.0$ against
$13.1\cdot10^{-3}$). The split is the same observed on the sine wave --- CfC
converges within one or two epochs, LTC and LRCU start slowly --- and it is to
be read with the same caution: after twenty epochs the second group has not yet
reached its own minimum, and a comparison at a fixed budget measures in part the
speed of convergence rather than the capacity. Separating the two would require
continuing the LTC and LRCU sessions until their validation loss turns, which
was not done. The skill on yaw follows the two groups too: between $-0.17$ and
$-0.24$ for the first, between $-1.0$ and $-4.4$ for the second.

\paragraph{The stateless control.}
The most relevant result of the comparison is where the control stands. The
\met{SparseConv_MLP}, identical to the others but with two linear layers in
place of the cells and therefore without memory from one step to the next, comes
second: its validation loss is $10\%$ higher than that of the
\met{SparseConv_CfC_NCP} and lower than that of the \met{SparseConv_CfC}, and
its per-component figures are in practice indistinguishable from those of the
two CfC networks. The difference between the best recurrent network and the
control, $0.46\cdot10^{-3}$, is smaller than the oscillation of the
\met{SparseConv_CfC_NCP} itself from one epoch to the next after its minimum
(between $4.8$ and $6.1\cdot10^{-3}$), and each configuration was run with a
single seed. Under the rule of Section~\ref{sec:exp-protocol}, then, the
advantage of the \met{SparseConv_CfC_NCP} over the control is indicative, not
established; what this means for the comparison as a whole is discussed in
Section~\ref{sec:res-discussion}.

\paragraph{Wiring.}
For the same cell, the NCP variant is better in all three cases: $4.63$ against
$5.52$ for CfC, $8.67$ against $11.78$ for LTC, $9.98$ against $13.03$ for LRCU.
For LTC and LRCU the difference, about a quarter of the loss, exceeds the
epoch-to-epoch oscillation; for CfC it is of the same order. The NCP wiring also
brings the two skip connections from the input to the second cell and from the
first cell to the output (Section~\ref{sec:lnn-wiring}), which shorten the
gradient path: with the data available the gain cannot be attributed to the
sparsity of the wiring rather than to the skips.

\paragraph{Cost.}
The cost per epoch varies little, between \num{557} seconds for the control and
\num{601} for the \met{SparseConv_LRCU_NCP}: on the corpus the time is dominated
by the encoder and by data loading, not by the cells. Nor does the parameter
count (Table~\ref{tab:architectures}) explain the ranking: the
\met{SparseConv_LRCU_NCP} is the largest network and comes fifth, the control is
the smallest and comes second.

\paragraph{The held-out test.}
The validation partition was used twice to decide --- to select the epoch and
to fit the calibration --- so its figures are not an unbiased estimate
(Section~\ref{sec:exp-protocol}). The test partition, which took part in
neither decision, was read once, on the seven selected checkpoints and after
the comparison above had been written: raw, with the calibration at the
identity, and calibrated, through the coefficients each checkpoint had fitted
on validation, which is the output the robot would receive.
Table~\ref{tab:res-test} reports it.

\begin{table}[htbp]
  \centering
  \small
  \caption{Corpus: the selected checkpoints on the held-out test partition
  (70 episodes, \num{17747} frames), read once. Losses in units of $10^{-3}$;
  the validation loss is repeated for reference. $R^2$ and skill on the raw
  output; skill on the forward component and on yaw. Zero predictor $127.7$,
  velocity echo $13.9$.}
  \label{tab:res-test}
  \setlength{\tabcolsep}{4pt}
  \begin{tabular}{@{}lrrrrrrrr@{}}
    \toprule
    \textbf{Network} & \textbf{val} & \textbf{test} & \textbf{cal.} & $R^2_{v_x}$ & $R^2_{v_z}$ & $R^2_{\omega_z}$ & \textbf{Skill} $v_x$ & \textbf{Skill} $\omega_z$ \\
    \midrule
    \met{CfC_NCP}     &  4.63 &  4.66 &  4.55 & 0.973 & 0.805 & 0.978 & 0.72 & $-0.16$ \\
    \met{MLP} (ctrl.) &  5.09 &  4.99 &  4.83 & 0.971 & 0.783 & 0.978 & 0.70 & $-0.17$ \\
    \met{CfC}         &  5.52 &  5.17 &  4.94 & 0.969 & 0.781 & 0.977 & 0.68 & $-0.20$ \\
    \met{LTC_NCP}     &  8.67 &  7.98 &  7.91 & 0.955 & 0.619 & 0.965 & 0.54 & $-0.84$ \\
    \met{LRCU_NCP}    &  9.98 &  9.39 &  9.34 & 0.948 & 0.575 & 0.952 & 0.47 & $-1.51$ \\
    \met{LTC}         & 11.78 & 12.08 & 12.06 & 0.939 & 0.513 & 0.912 & 0.37 & $-3.61$ \\
    \met{LRCU}        & 13.03 & 12.55 & 12.45 & 0.942 & 0.443 & 0.908 & 0.40 & $-3.83$ \\
    \bottomrule
  \end{tabular}
\end{table}

Three things follow from it. The first is that the validation figures were not
flattered by being used to decide: for every network the test loss lies within
8\% of the validation loss, in both directions, and the ranking of the seven is
the same on the two partitions. The second is that the calibration fitted on
validation carries over to data it has never seen: the calibrated test loss is
lower than the raw one for all seven networks, by up to 4\% for the
\met{SparseConv_CfC}, and on the lateral component it again brings a negative
raw $R^2$ back to zero. The third is that none of the readings of this section
is an artefact of the validation partition. The best recurrent network and the
control remain separated by less than the epoch-to-epoch fluctuation ($4.66$
against $4.99$); no network beats the echo on yaw; the restart ratio is, if
anything, lower still, at most $0.07$ and slightly negative for the control; the
error on still frames exceeds the error on moving ones by a factor between
$1.6$ and $3.0$; and the forward $R^2$ stays between $0.93$ and $0.98$ on both
platforms.


\subsection{What the Encoder Learns}
\label{sec:res-sl-encoder}

The sections above read the command; this one reads the encoder that feeds it,
on the validation partition and for all seven networks. Three measurements are
available, and they answer questions of increasing strictness. The first is the
auxiliary head itself (Section~\ref{sec:geometry-head}): at the selected epoch
it explains about $47$--$50\%$ of the variance of the per-sector clearances and
$69$--$81\%$ of that of the height. That figure is an upper reading of what the
encoder represents, because the head is non-linear and trained together with
the encoder. The two tools of Section~\ref{sec:exp-tooling} ask the stricter
questions: whether the command actually \emph{uses} what the encoder computes ---
the ablation --- and whether a \emph{linear} read-out of the frozen descriptor,
fitted on some episodes and scored on others, recovers the geometry --- the
probe. The probe's figures are therefore lower than the head's, and they are the
ones to read.

\paragraph{Does the command use the map?}
The ablation replaces, with the same weights, the encoder's per-sample output
by its batch mean --- the information is gone, the fusion stays in the range it
was trained on --- and measures the calibrated validation loss again
(Table~\ref{tab:res-encoder}). The loss rises for all seven networks, from
$+8\%$ for the \met{SparseConv_LRCU} to $+22\%$ for the two CfC networks, and
it rises further, by $14$--$48\%$, when every observation is given another
sample's cloud instead. The command therefore depends on the content of the
cloud and not merely on its presence as a constant, and the better a network
imitates, the more it relies on it. The dependence is at least as strong on the
$143^3$ grid, the one the fleet flies ($+7$ to $+25\%$). Per component, what the
map carries is mostly the vertical command: for the \met{SparseConv_CfC_NCP}
its $R^2$ falls from $0.77$ to $0.72$ without the encoder, the forward one from
$0.975$ to $0.968$, and the yaw one does not move. The finest grid, $363^3$,
holds only UGV frames with a residual loss four times smaller than the others,
and there the percentages scatter between $-35\%$ and $+146\%$; it is not read.

\paragraph{Does the encoder represent the scene?}
The probe answers the second question independently of the command: a ridge
regression from the frozen descriptor to the geometry each frame's own cloud
shows --- the clearance per azimuth sector and the height above the ground ---
fitted on half of the validation episodes and scored on the other half, so that
a descriptor identifying the episode rather than describing the scene gains
nothing. The same architecture at its initialisation gives the floor a trained
read-out must clear (Table~\ref{tab:res-encoder}). Untrained, the descriptor
already explains $24\%$ of the variance of the sector clearances and $43\%$ of
that of the height, which is what a random sparse convolution preserves of a
point cloud. Trained, the seven encoders reach $37$--$44\%$ on the sectors and
$68$--$75\%$ on the height, with differences between architectures that are
small and do not follow their ranking in the loss: the encoder is the same for
all of them, and what the cells do downstream does not change what it learns.
The height is the less telling of the two figures, since it separates the two
platforms before anything else --- the UGV sits at a fixed height, the UAV
cruises --- and a descriptor that recognises the robot explains most of it.

The sector clearances are the figure that matters for the fleet, and on the
$143^3$ grid the fleet flies they are the weakest: between $0.15$ and $0.29$
against an untrained floor of $0.15$, so that for some networks the directional
distance to the obstacles is barely more readable after training than before.
Training does make the descriptor more informative in the other two respects the
probe measures: its spread across samples grows from $0.07$ to $0.16$--$0.22$ on
that grid, and its displacement under a shift of three voxels falls from $0.80$
to $0.44$--$0.58$ of the distance between two different clouds. Taken together
with the ablation, the picture is consistent: the imitation uses the map, it
uses it mostly for the vertical command, and it teaches the encoder where the
obstacles are only in part --- the part the reinforcement stage has to build on.

\begin{table}[htbp]
  \centering
  \small
  \caption{Corpus: what the command takes from the encoder and what the
  encoder represents, on the validation partition. \emph{Ablation}: increase of
  the calibrated loss when the encoder's output is replaced by its batch mean,
  overall and on the $143^3$ grid the fleet flies, and when the clouds are
  shuffled across the batch. \emph{Probe}: fraction of the variance of the
  sector clearances and of the height explained by a linear read-out of the
  frozen descriptor, on held-out episodes; the last row is the same
  architecture at its initialisation.}
  \label{tab:res-encoder}
  \setlength{\tabcolsep}{4pt}
  \begin{tabular}{@{}lrrrrrr@{}}
    \toprule
    & \multicolumn{3}{c}{\textbf{Ablation}} & \multicolumn{3}{c}{\textbf{Probe}} \\
    \cmidrule(lr){2-4}\cmidrule(l){5-7}
    \textbf{Network} & \textbf{Mean} & \textbf{Mean, $143^3$} & \textbf{Shuffled} & \textbf{Sectors} & \textbf{Sectors, $143^3$} & \textbf{Height} \\
    \midrule
    \met{CfC_NCP}     & $+21.9\%$ & $+24.9\%$ & $+40.5\%$ & 0.40 & 0.18 & 0.70 \\
    \met{MLP} (ctrl.) & $+18.8\%$ & $+20.4\%$ & $+47.9\%$ & 0.40 & 0.19 & 0.73 \\
    \met{CfC}         & $+22.4\%$ & $+22.5\%$ & $+39.0\%$ & 0.42 & 0.29 & 0.71 \\
    \met{LTC_NCP}     & $+15.3\%$ & $+19.6\%$ & $+32.1\%$ & 0.44 & 0.15 & 0.73 \\
    \met{LRCU_NCP}    & $+13.7\%$ & $+21.1\%$ & $+24.3\%$ & 0.37 & 0.21 & 0.75 \\
    \met{LTC}         & $+11.2\%$ & $+13.7\%$ & $+17.6\%$ & 0.44 & 0.20 & 0.72 \\
    \met{LRCU}        &  $+8.2\%$ &  $+7.3\%$ & $+14.3\%$ & 0.43 & 0.22 & 0.68 \\
    \midrule
    untrained         & -- & -- & -- & 0.24 & 0.15 & 0.43 \\
    \bottomrule
  \end{tabular}
\end{table}


\section{Reinforcement Learning}
\label{sec:res-rl}

\subsection{Validation on LunarLander}
\label{sec:res-rl-lunar}

\paragraph{Setting.}
The recurrent agent was verified on the LunarLander environment of
\texttt{gymnasium} (version~3), with a dense eight-component observation and
four discrete actions, without the fleet. The seven vector networks are trained
from scratch on \num{16} parallel environments, with rollouts of \num{512} steps
per environment, that is \num{8192} transitions per iteration, for \num{733}
iterations and about \num{6.0} million steps. The six recurrent networks share
a single configuration (Section~\ref{sec:rl-presets}): windows of \num{32}
steps, six epochs per update, a target KL of $0.02$, a clipping $\epsilon$
annealed from $0.2$ to $0.08$ and a learning rate going from $2\cdot10^{-4}$ to
$6\cdot10^{-5}$. The control has its own configuration but the same step
budget, so that the curves are comparable along the horizontal axis. Each
network was run once. The external reference is the conventional threshold of
$200$ mean return, above which LunarLander is considered solved; times to
threshold are read on a moving average over twenty iterations.

\begin{table}[htbp]
  \centering
  \small
  \caption{LunarLander: steps to threshold (millions, moving average over 20
  iterations), mean return, episode length and worst episodic return over the
  last 50 iterations, wall-clock duration of the session, and largest recorded
  norm of the recurrent state.}
  \label{tab:res-lunar}
  \setlength{\tabcolsep}{4pt}
  \begin{tabular}{@{}lrrrrrrr@{}}
    \toprule
    \textbf{Network} & $\geq 200$ & $\geq 250$ & \textbf{Return} & \textbf{Length} & \textbf{Worst} & \textbf{Time} & $\lVert h\rVert_{\max}$ \\
    \midrule
    \met{MLP} (ctrl.) & 1.02 & 1.34 & 284.3 & 162 &   14 & 22 min & -- \\
    \met{LRCU_NCP}     & 1.01 & 1.43 & 276.4 & 193 & $-314$ & 90 min & 299 \\
    \met{CfC}          & 1.54 & 2.37 & 284.9 & 172 &    0 & 45 min & 5.9 \\
    \met{LRCU}         & 1.65 & 2.25 & 276.7 & 186 & $-144$ & 73 min & 392 \\
    \met{CfC_NCP}      & 1.83 & 2.29 & 272.6 & 286 &   19 & 49 min & 5.2 \\
    \met{LTC}          & 2.16 & 3.13 & 275.3 & 193 &  $-22$ & 54 min & 3.2 \\
    \met{LTC_NCP}      & 2.08 & 3.71 & 268.0 & 221 &  $-40$ & 63 min & 3.1 \\
    \bottomrule
  \end{tabular}
\end{table}

\paragraph{Outcome.}
All seven networks solve the task (Table~\ref{tab:res-lunar},
Figure~\ref{fig:res-lunar}). The threshold of $200$ is crossed between one and
slightly more than two million steps, and once crossed the moving average does
not fall back below it by more than a few points ($195$ at the lowest, for the
\met{LNN_LRCU}). Over the last fifty iterations the mean return lies between
$268$ and $285$, a range of seventeen points against a spread between the
episodes of a single iteration of $30$--$47$ points. With a single run per
network, the networks are therefore to be considered equivalent on the final
result; the readable differences are in speed and in cost.

The course is the same for every network: in the first \num{500000}--\num{800000}
steps the episode length rises to the environment's limit of a thousand steps
--- the policy has learnt not to crash and hovers until the time runs out ---
and only afterwards does it learn to land, the length dropping to
$160$--$220$ steps as the return crosses $200$. It is the known sequence for this
environment, and observing it is itself a check on collection, advantage
estimation and update.

\begin{figure}[htbp]
  \centering
  \includegraphics[width=\textwidth]{ch08/lunar_performance.png}
  \caption{LunarLander: episodic return (mean, standard deviation, minimum and
  maximum), episode length, episodes per iteration, and raw and smoothed
  return, for the seven vector networks.}
  \label{fig:res-lunar}
\end{figure}

\paragraph{Correctness of the recurrent machinery.}
The task verifies above all what a return curve does not show. The
\met{replay_logratio_max} indicator, which compares the action probabilities
recomputed by replaying the recurrent state with those recorded during
collection and must be zero up to numerical error, never exceeds
$3\cdot10^{-5}$ in any iteration of any session: the state is reconstructed
correctly at update time. The explained variance of the critic ends between
$0.93$ and $0.97$; the early stop on the KL divergence fires in fewer than 2\%
of the iterations, and never in the last fifty, because the policy step has
shrunk to a divergence of the order of $10^{-3}$. The entropy falls from
$\ln 4 \approx 1.39$, the uniform distribution over the four actions, to
$0.34$--$0.55$.

\paragraph{Speed and cost.}
The control reaches $250$ first and ties with the best network on the final
return. This is expected: the LunarLander observation contains the full state of
the vehicle, so the test verifies the \emph{machinery} --- that the recurrent
networks learn as much as a network with no need to remember --- and not the
value of memory. Among the recurrent networks the \met{LNN_LRCU_NCP} is the
fastest to reach $200$, level with the control, and the two LTC networks are the
slowest, the \met{LNN_LTC_NCP} reaching $250$ only after $3.7$ million steps.
In wall-clock time the order reverses, because here there is no encoder to
dominate the cost and the cells weigh: for the same number of steps the control
takes $22$ minutes, the CfC networks $45$--$49$, the LTC networks $54$--$63$ and
the LRCU networks $73$--$90$. The \met{LNN_LRCU_NCP} is thus the most
sample-efficient network and the least time-efficient one. The only late
anomaly is the \met{LNN_CfC_NCP}'s: between $4.2$ and $5.8$ million steps part
of its episodes go back to hovering before landing, and the mean length over
the last fifty iterations ($286$ steps) still bears its trace.

\paragraph{State norms.}
The sharpest difference between the cells is not in the return but in the state
(Figure~\ref{fig:res-lunar-states}). For CfC and LTC the largest state norm
stays below $6$ throughout the session; for the two LRCU networks it reaches
$300$--$390$. It is the growth already glimpsed on the sine wave, amplified by
episodes twenty times longer. It does not prevent learning --- the two LRCU
networks end at $276$ --- but it coincides with the only two cases of disastrous
episodes at the end of the session (worst return $-314$ and $-144$ over the last
fifty iterations, against non-negative values for CfC and the control), and it
is a signal to be watched on the fleet, where episodes are longer and the time
interval is not constant. The stability of LTC is instead the result of a
correction made on this very environment: with the hyperbolic tangent as
activation the fused step of Equation~\eqref{eq:ltc-fused} is not stable for
negative values of $f$, and an earlier session had driven the state norm to
$3\cdot10^{11}$ within \num{41} iterations. With the sigmoid, which is the original formulation
(Section~\ref{sec:cell-ltc}), $f$ is non-negative and the state stays bounded.

\begin{figure}[htbp]
  \centering
  \includegraphics[width=\textwidth]{ch08/lunar_states.png}
  \caption{LunarLander: recurrent state norms per cell, for the six recurrent
  networks.}
  \label{fig:res-lunar-states}
\end{figure}

\paragraph{What the test does not verify.}
The action space is discrete, so the Gaussian policy, its learnt width and the
envelope penalty (Section~\ref{sec:rl-actions}) are not exercised; the time step
is constant; and the networks start from scratch, so the critic warm-up is not
exercised either (Section~\ref{sec:exp-protocol}). The test establishes that
the recurrent algorithm is correct and that each cell learns within it.


\subsection{Training on the Fleet}
\label{sec:res-rl-fleet}
% Sources: ~/WorkEnv/ai/tests/rl/cur_*_01/training_log.csv (7 runs, 320
% iterations each); figures from ~/WorkEnv/ai/plots/rl/cur_all/ copied to
% images/ch08/fleet_* (fleet_outcomes.png regenerated for the thesis).
% "Imitated" = iterations 1-19 (critic warm-up, actor frozen); "final" = last
% 50 iterations (271-320). Shares are per episode, each iteration weighted by
% the number of episodes that closed in it.

\paragraph{Setting.}
The seven spatial networks were fine-tuned on the fleet, each starting from its
own imitation checkpoint, the one selected in
Section~\ref{sec:res-sl-comparison}. Every session lasts \num{320} iterations
of \num{2048} transitions, that is \num{655360} environment steps, about a
ninth of the LunarLander budget; the first \num{19} iterations are reserved for
the critic (Section~\ref{sec:rl-finetuning}), and for their whole duration the
actor stays frozen bit for bit.

\paragraph{A fleet of graded difficulty.}
The fleet counts eight slices, the number the workstation of
Section~\ref{sec:exp-tooling} holds beside the training session: the slices
and the session share its \SI{16}{\gibi\byte} of accelerator memory, and with
a trainable encoder the session needs a large share of it, so that the size of
the fleet was set by the machine rather than by the throughput. The eight are
four multirotors and four rovers, all in the
\texttt{warehouse\_small} world --- a $40\times24$\,m warehouse --- and all on
the $143^3$ grid, with \SI{0.1}{\metre} voxels and a map radius of
\SI{7}{\metre}: the same grid on which Section~\ref{sec:res-sl-encoder} measures
the ablation and the probe. The slices do not pose the same task. Each platform
has four levels of difficulty, each assigned to one slice for the whole session
(Table~\ref{tab:res-fleet-slices}): a short trip with the goal ahead, one of the
same length with the goal to the side, one of \SI{5.5}{\metre} and one of
\SI{7}{\metre}, on the last two of which the straight route is clear only in
part of the episodes. This is not a curriculum: the levels do not follow one
another, they all run together from the first step to the last, and the trainer
knows nothing of them. Every batch of data therefore holds easy and hard trips
in fixed proportions, on both platforms. Spawn and goal are fixed by the
configuration of each slice and are not drawn at random: every slice poses the
same trip at every episode.

\begin{table}[htbp]
  \centering
  \small
  \caption{Fleet: the eight slices, all in the \texttt{warehouse\_small} world
  on the $143^3$ grid. \emph{Clear route} is the share of episodes in which the
  straight segment from spawn to goal crosses no obstacle; spawn and goal are
  fixed per slice. \emph{Budget} is the maximum number of steps of an episode,
  and $\gamma$ the discount derived from it.}
  \label{tab:res-fleet-slices}
  \begin{tabular}{@{}clcllrr@{}}
    \toprule
    \textbf{Slice} & \textbf{Robot} & \textbf{Level} & \textbf{Trip} & \textbf{Clear route} & \textbf{Budget} & $\gamma$ \\
    \midrule
    0 & UAV & 1 & \SI{2.5}{\metre}, ahead & always & 94  & 0.9894 \\
    1 & UAV & 2 & \SI{2.5}{\metre}, at $90^\circ$ & always & 94  & 0.9894 \\
    2 & UAV & 3 & \SI{5.5}{\metre} & $45\%$ & 207 & 0.9952 \\
    3 & UAV & 4 & \SI{7}{\metre}   & $8\%$  & 263 & 0.9962 \\
    4 & UGV & 1 & \SI{2.5}{\metre}, ahead & always & 95  & 0.9895 \\
    5 & UGV & 2 & \SI{2.5}{\metre}, at $90^\circ$ & always & 95  & 0.9895 \\
    6 & UGV & 3 & \SI{5.5}{\metre} & $43\%$ & 207 & 0.9952 \\
    7 & UGV & 4 & \SI{7}{\metre}   & $0\%$  & 263 & 0.9962 \\
    \bottomrule
  \end{tabular}
\end{table}

\paragraph{Episodes, reward and optimisation.}
The budget of every episode is three times the ideal time at maximum speed
(Section~\ref{sec:rl-episodes}), and the discount is derived per slice,
$\gamma = 1 - 1/T$ (Section~\ref{sec:rl-presets}); an episode arrives within
\SI{1.0}{\metre} of the goal, ten voxels. No collision margin is armed; the
stall condition, armed on the rovers only, ends the episode after twelve
consecutive steps in which the commanded vehicle does not move. The reward is
the \texttt{safe\_progress} policy of Section~\ref{sec:rl-reward} with its
default values --- an arrival bonus of $25$, a penalty of $10$ for losing the
vehicle --- and a cruise altitude of \SI{2.2}{\metre} for the multirotors. The
encoder is trainable, without the auxiliary geometry term.

All seven networks use the same hyper-parameters: the two presets they were
taken from are identical, and when a session starts from a checkpoint the
network is rebuilt from its metadata, so that the preset fixes only the
optimisation. The learning rate is $3\cdot10^{-5}$ for the actor and
$3\cdot10^{-4}$ for the critic, the entropy bonus $0.001$, the target KL
$0.02$, with windows of eight steps and four epochs per update. These values
are meant to fine-tune without drifting away from the imitated policy, and one
effect is visible in the logs: the width of the Gaussian, which starts at 10\%
of the envelope, is trainable but moves by at most 3\% in \num{320}
iterations, and its mean stays at $0.080$ for the whole session in every
network.

Each network was trained once, with seed $0$. A session closes between
\num{5507} and \num{8048} episodes and costs about fifteen hours, at about
\num{12} steps per second for the whole fleet: a step costs a hundred to four
hundred times as much as on LunarLander, because it is a physics simulation, a
lidar rendering and a full pass through the pipeline of
Chapter~\ref{ch:pipeline}. The cost is dominated by the environment and does
not depend on the cell, and the budget of \num{320} iterations is therefore a
time constraint rather than a choice of convergence.

Consistently with Section~\ref{sec:metrics-rl}, the result is read on the share
of episodes that reach the goal and on the reasons episodes end, not on the
return: the shares are computed per episode, each iteration weighted by the
number of episodes that closed in it.

\paragraph{The imitated policy on the fleet.}
The critic warm-up iterations are the only free measurement of the imitated
policy on the fleet, and it is from them --- not from zero --- that the gain of
the fine-tuning is to be read (Section~\ref{sec:exp-protocol}). The measurement
is clear-cut: \emph{no} imitated policy flies the task. The arrival rate goes
from 0.6\% for the control to 9.4\% for the \met{SparseConv_CfC_NCP}, and the
most frequent way for an episode to end is to run out of time: between 45\%
and 76\% of the episodes end on \met{max_steps}, the rest splitting between
tipping over (9--27\%) and stalling (4--14\%). The mean progress per episode
is negative for every network, between $-0.3$ and $-1.1$\,m: on average the
vehicle ends the episode further from the goal than it started. It is the
behaviour the restart metrics of the corpus anticipated
(Section~\ref{sec:res-sl-corpus}): a policy that at departure commands 2\% of
what the pilot commanded does not depart, and uses up the time it is given.

\paragraph{Outcome.}
After \num{320} iterations all seven networks reach the goal in between 36\%
and 63\% of the episodes (Table~\ref{tab:res-fleet},
Figure~\ref{fig:res-fleet-outcomes}), that is five to a hundred times as often
as the imitated policy they started from. The mean return per episode rises
from between $-5.7$ and $-1.2$ to between $6.3$ and $10.9$, the mean episode
length falls from \num{99}--\num{121} to \num{66}--\num{112} steps, and the
mean progress turns positive, between $0.5$ and $1.2$\,m per episode.

The course is not gradual. For roughly the first hundred iterations the arrival
rate stays below 20\% for every network, and running out of time remains the
dominant end reason; between iterations $100$ and $220$ the rate rises, in
every network, at the same moment the progress term of the reward turns
positive, that is when the policy starts closing distance instead of losing it.
Afterwards the growth slows down, and at the end of the session three networks
--- \met{SparseConv_CfC}, the control and \met{SparseConv_LTC_NCP} --- still
have their moving average at its maximum: the step budget stopped the training
before the curve flattened, and two more networks reach their maximum within
the last twenty iterations. For the remaining two the best window came earlier
--- 61\% at iteration \num{219} for the \met{SparseConv_CfC_NCP} and 48\% at
iteration \num{281} for the \met{SparseConv_LRCU} --- but the comparison is read
on the final iterations for every network, as Section~\ref{sec:exp-protocol}
prescribes, and the choice of the checkpoint to deploy is left to a later
stage.

\begin{table}[htbp]
  \centering
  \small
  \caption{Fleet: arrival rate of the imitated policy (critic warm-up
  iterations) and after fine-tuning (last 50 iterations, and separately in its
  two halves of 25), end reasons over the last 50 iterations, mean return per
  episode, and largest norm of the recurrent state.}
  \label{tab:res-fleet}
  \setlength{\tabcolsep}{3.5pt}
  \begin{tabular}{@{}lrrrrrrrr@{}}
    \toprule
    & \multicolumn{3}{c}{\textbf{Arrival rate}} & \multicolumn{3}{c}{\textbf{End reason (final)}} & & \\
    \cmidrule(lr){2-4}\cmidrule(lr){5-7}
    \textbf{Network} & \textbf{Imitated} & \textbf{Final} & \textbf{Halves} & \textbf{Flipped} & \textbf{Time} & \textbf{Stalled} & \textbf{Return} & $\lVert h\rVert_{\max}$ \\
    \midrule
    \met{CfC}         & 0.078 & 0.627 & 0.61 / 0.64 & 0.09 & 0.27 & 0.02 & 10.8 & 5.0 \\
    \met{MLP} (ctrl.) & 0.006 & 0.611 & 0.57 / 0.65 & 0.11 & 0.20 & 0.08 & 10.9 & -- \\
    \met{LTC_NCP}     & 0.038 & 0.471 & 0.45 / 0.49 & 0.04 & 0.35 & 0.13 &  7.4 & 3.2 \\
    \met{CfC_NCP}     & 0.094 & 0.464 & 0.46 / 0.47 & 0.03 & 0.38 & 0.13 &  6.5 & 2.9 \\
    \met{LRCU}        & 0.069 & 0.405 & 0.43 / 0.38 & 0.06 & 0.50 & 0.03 &  6.3 & 76.8 \\
    \met{LRCU_NCP}    & 0.021 & 0.385 & 0.36 / 0.41 & 0.10 & 0.50 & 0.02 &  6.9 & 6.9 \\
    \met{LTC}         & 0.051 & 0.356 & 0.35 / 0.36 & 0.06 & 0.54 & 0.05 &  7.5 & 3.6 \\
    \bottomrule
  \end{tabular}
\end{table}

\begin{figure}[htbp]
  \centering
  \includegraphics[width=\textwidth]{ch08/fleet_outcomes.png}
  \caption{Fleet: how episodes end over the \num{320} iterations, one panel per
  network. Shares per episode, as a moving average over 20 iterations weighted
  by the number of episodes; arrivals sit at the baseline so that they compare
  across panels. The grey band is the critic warm-up, that is the imitated
  policy; the figure in each panel is the arrival rate over the last 50
  iterations (Table~\ref{tab:res-fleet}).}
  \label{fig:res-fleet-outcomes}
\end{figure}

\paragraph{How episodes end.}
Most of the gain comes from episodes that used to end on the clock: for six
networks out of seven the share of \met{max_steps} falls by $16$--$49$ points;
for the \met{SparseConv_LRCU}, which started from the lowest share (45\%), it
rises slightly, to 50\%. Tipping over is roughly halved in every network, from
9--27\% to 3--11\%, whereas stalling moves in different directions: it falls
for four networks, stays level for the control, and rises for the two NCP
networks with the CfC and LTC cells, from 4\% to 13\% and from 11\% to 13\%.
Fine-tuning, in other words, taught above all how to \emph{depart} and arrive
in time, and secondarily how to lose the vehicle less often. Even at the end of
the session, running out of time remains the main cause of failure for every
network, between 20\% and 54\% of the episodes.

\paragraph{What the return is made of.}
The decomposition of the reward (Section~\ref{sec:metrics-rl-reward}) shows
that, at the end of the session, the mean reward per step coincides in practice
with the terminal term alone: the ratio between the two lies between $0.96$
and $1.08$ for every network. The progress earned ($+0.005$ to $+0.018$ per
step) and the dense costs --- time, fixed at $-0.0064$, safety ($-0.003$ to
$-0.005$) and altitude ($-0.0005$ to $-0.0014$) --- cancel almost exactly, and
the return is in practice a count of arrivals. This is consistent with what the
reward budget prescribes (Section~\ref{sec:rl-reward}): the dense costs must not
be able to cost more than progress gives back, and they do not. It follows,
however, that the return adds no information to the arrival rate, which is why
Table~\ref{tab:res-fleet} reports it for completeness only.

\paragraph{Optimisation.}
The critic starts from nothing and it shows: during the warm-up the explained
variance lies between $0.17$ and $0.36$, at the end of the session between
$0.88$ and $0.94$. The trust region is instead the binding constraint for
almost the whole session: after the warm-up the actor's update is stopped early
in 58--79\% of the iterations for six networks out of seven, and in 31\% for
the \met{SparseConv_LTC}. Every iteration therefore moves the policy as far as
the region allows and no further, and it is plausibly this, together with the
step budget, that sets the pace at which the arrival rate grows.

\begin{figure}[htbp]
  \centering
  \includegraphics[width=\textwidth]{ch08/fleet_performance.png}
  \caption{Fleet: episodic return (mean, standard deviation, minimum and
  maximum), episode length, episodes per iteration, and raw and smoothed
  return, for the seven spatial networks.}
  \label{fig:res-fleet-performance}
\end{figure}

\paragraph{The preliminary sessions.}
The configuration used is the last of a series, and four measurements from the
earlier sessions explain why the result is read on the end reasons and on the
progress rather than on the return. In one session the share of episodes
ending tipped over went from 73\% to 97\% while the return improved from $-19$
to $-4.6$: the policy had learnt to lose the vehicle sooner, paying fewer dense
costs. In another, of \num{650} iterations, the mean progress per episode fell
from $1.26$ to $0.04$\,m while the return rose throughout. In both cases the
return improved while the task got worse. The other two measurements concern
the reward itself: the progress term in its canonical form, telescoped over an
episode, is worth $d(s_0)(1-\gamma^T)$, that is $+5.6$ on a slice with a short
trip and $+24.5$ on the one with the long trip, against an arrival bonus of
$20$ --- on that slice, standing still paid more than arriving --- and in a
later configuration the terminal term alone carried 90\% of the return. These
measurements are what led to the undiscounted progress term and to the double
budget constraint of Section~\ref{sec:rl-reward}.


\subsection{Architecture Comparison}
\label{sec:res-rl-comparison}

The conditions of the comparison are those of Section~\ref{sec:exp-protocol}:
the same step budget, the same fleet, the same reward, and for every network the
imitation checkpoint of the same architecture. The optimisation hyper-parameters
are also the same for all the networks.

\paragraph{Two groups, again.}
The networks split into two groups. The \met{SparseConv_CfC} and the control
arrive in 61--63\% of the episodes; the other five between 36\% and 47\%. The
separation between the two groups, about $15$ points, exceeds the oscillation
within each run: between the first and the second half of the last fifty
iterations the arrival rate of the same network changes by at most $8$ points.
Within each group, by contrast, the differences are of the same order as that
oscillation and are not to be read as a ranking. With a single run per network,
and with a fleet whose interleaving across slices depends on the load of the
machine (Section~\ref{sec:exp-protocol}), even the separation between the two
groups remains indicative rather than established.

\paragraph{The stateless control.}
The network without memory arrives as often as the best recurrent one, although
it started from the worst imitated policy of all (0.6\%), and although on the
fleet the task is partially observable: the cloud sees only within the map
radius and the velocity is estimated. What this implies for the comparison is
discussed in Section~\ref{sec:res-discussion}.

\paragraph{Wiring and cells.}
The NCP wiring, which on the corpus improved the loss for the same cell, brings
no consistent advantage on the fleet: the \met{SparseConv_CfC} beats its NCP
variant ($0.63$ against $0.46$), the \met{SparseConv_LTC_NCP} beats the
\met{SparseConv_LTC} ($0.47$ against $0.36$), and the two LRCU networks are
equivalent ($0.41$ and $0.39$). The \met{SparseConv_CfC_NCP} also shows the
largest late regression: its moving average reaches 61\% at iteration \num{219}
and falls back to 46\% by the end, while the episode length and the number of
episodes per iteration follow the same course
(Figure~\ref{fig:res-fleet-performance}). Without a second run it cannot be
said whether this is an instability of the network or an oscillation of the
session.

The LRCU cells confirm on the fleet what the sine wave and LunarLander had
shown. The largest state norm of the \met{SparseConv_LRCU} is $72$--$77$
throughout the session, against $3$--$7$ for all the others
(Figure~\ref{fig:res-fleet-states}); the NCP variant stays below $7$. It is the
only network whose arrival rate drops between the two final halves (43\% and
38\%). The LRCU networks are also those whose imitated policy most often
commands outside the envelope: during the warm-up the mean of the Gaussian of
the \met{SparseConv_LRCU_NCP} falls outside the limits in 16.5\% of the steps
and that of the \met{SparseConv_LRCU} in 13.2\%, against less than 1\% for the
others; fine-tuning brings both below 5\%.

\begin{figure}[htbp]
  \centering
  \includegraphics[width=\textwidth]{ch08/fleet_states.png}
  \caption{Fleet: recurrent state norms per cell, for the six recurrent
  networks.}
  \label{fig:res-fleet-states}
\end{figure}


\section{Cross-Stage Analysis}
\label{sec:res-cross}

The two stages were evaluated separately; this section asks what the first says
about the second. The sample is seven networks, one run each: with seven points
a rank correlation is significant at the 5\% level only beyond
$|\rho| \approx 0.79$, and what follows is to be read against that threshold.

\paragraph{The imitation loss does not predict how the imitated policy flies.}
The ranking of the networks by test loss (Table~\ref{tab:res-test}) and the
ranking by arrival rate of the imitated policy on the fleet are unrelated
($\rho = -0.14$). The extreme case is the control: second by loss, last by
arrivals (0.6\%). This is not surprising: the mean loss is dominated by the
moving frames and by the forward velocity, whereas what decides whether an
episode arrives is the departure --- the rare event on which every imitated
network fails in the same way (Section~\ref{sec:res-sl-corpus}). The loss
measures well what it measures, and that is not what the fleet needs.

\paragraph{The quality of the imitation weakly predicts the end point.}
The ranking by loss is, on the other hand, related to the \emph{final} arrival
rate ($\rho = -0.64$) and to the gain of the fine-tuning ($\rho = -0.75$): the
three networks with the lowest loss are three of the four with the best final
result, and the CfC-plus-control group is the same in the two stages. Neither
coefficient clears the threshold, and the two groups also coincide with the
speed of convergence observed on the corpus, so that the quality of the
starting point cannot be separated from how easily the network optimises.

\paragraph{How much the network uses the map predicts the gain.}
The strongest relation is with the encoder ablation (Table~\ref{tab:res-encoder}):
how much the validation loss degrades when the information of the cloud is
removed is related to the final arrival rate and to the gain of the fine-tuning
($\rho = +0.79$ for both, at the threshold). The networks whose imitated command
depends most on the map are those that fine-tuning takes furthest. The probe,
by contrast, is related to neither ($|\rho| \le 0.32$, on the $143^3$ grid
alone as well): that the descriptor \emph{contains} the geometry predicts
nothing; that the command \emph{uses} it does. This is the most useful
indication for future work, because it suggests selecting the imitation
checkpoint on its dependence on the map as well, and not on the loss alone.

\paragraph{What fine-tuning corrected.}
Read together with Section~\ref{sec:res-sl-corpus}, the result on the fleet
matches point by point the main defect of the imitation. The corpus showed a
policy that at a restart commands 2\% of what the pilot commanded; the fleet
shows imitated policies that run out of time in 45--76\% of the episodes with
negative progress; fine-tuning removes up to fifty points from that end reason
and moves almost all of them to arrivals. What fine-tuning corrected only in
part is just as readable: tipping over is halved, but stalls --- contact with
an obstacle for the vehicles that do not fly --- do not fall consistently, and
they rise precisely for two of the networks that arrive best. This is
consistent with what the probe found on the fleet's grid: the directional
distance to the obstacles is the part of the scene the encoder represents
worst.


\section{Discussion, Limitations and Threats to Validity}
\label{sec:res-discussion}

\paragraph{What the results establish.}
The whole chain works: the three cells learn a known dynamics, the recurrent
agent is correct and learns on a reference problem, the imitation of the corpus
generalises to unseen episodes with an error a third of that of echoing the
velocity, and fine-tuning on the fleet takes the arrival rate from below 10\% to
between 36\% and 63\% for every architecture. The two-stage procedure therefore
produces a policy that flies a task imitation alone does not.

\paragraph{What they do not establish.}
The comparison between architectures does not single out a winner, and the
reason is the control. In all three stages the network without state is among
the best: second by loss on the corpus, the fastest on LunarLander, level with
the best recurrent network on the fleet. In each case the difference between
the control and the best recurrent network is smaller than the variability
within a single run, so that on the tasks used memory has not shown that it is
worth its state.

The three stages say this for different reasons. On LunarLander it is expected,
because the observation contains the full state of the vehicle. On the corpus
the observation already carries the robot's velocity, the position of the goal
and, through the encoder, the local geometry, so that the part of the command
imitation can learn is largely a function of the current state; and where the
demonstration is ambiguous --- the restart --- no network, with or without state,
reaches $10\%$ of what the pilot commanded. On the fleet the task is partially
observable, and the control still arrives as often as the best recurrent
network. The hypothesis that continuous-time cells bring an advantage on this
problem is therefore not refuted, but it is not supported by the data either.

\paragraph{Threats to internal validity.}
\begin{itemize}
  \item \emph{One run per configuration.} Each network was trained once per
        stage. The rule of Section~\ref{sec:exp-protocol} asks for the spread
        across seeds to be measured before a difference is stated; here it was
        replaced by the oscillation within a run (across epochs, across halves of
        a window), which underestimates it. All differences between
        architectures are therefore reported as indicative.
  \item \emph{Fixed budgets.} Twenty epochs on the corpus and \num{320}
        iterations on the fleet are time constraints. On the corpus LTC and LRCU
        were still descending at the last epoch; on the fleet three networks were
        still at their maximum. The comparison measures in part the speed of
        convergence.
  \item \emph{Asynchrony of the fleet.} Two identical runs do not collect the
        same sequence of transitions; the seed fixes the tasks and not their
        order.
  \item \emph{Shared hyper-parameters.} The networks use the same optimisation
        configuration, tuned on the \met{SparseConv_CfC_NCP}; it may suit some
        cells better than others.
  \item \emph{Fleet-level metrics.} Arrival rates and end reasons are recorded
        over the whole fleet, not per slice: multirotors and rovers cannot be
        separated, nor can the four levels. A difference between architectures
        concentrated on the hard trips may therefore be diluted by the easy
        ones.
\end{itemize}

\paragraph{Threats to external validity.}
\begin{itemize}
  \item \emph{Simulation only.} No policy has been run on a real robot. The
        observation pipeline is the same as on board
        (Section~\ref{sec:rl-environment}), but the dynamics, the sensors and
        the latencies are those of the simulator.
  \item \emph{One world, four fixed levels.} Fine-tuning took place in a single
        world and on trips between \SI{2.5}{\metre} and \SI{7}{\metre}; nothing
        says how the policy behaves on more cluttered maps or on longer trips.
  \item \emph{A small, unbalanced corpus.} \num{476} episodes, two robots, two
        worlds; 86\% of the energy of the command on a single component, the
        lateral velocity absent. The probe finds the directional geometry poorly
        readable precisely on the fleet's grid, which limits what imitation can
        transfer.
  \item \emph{Validation used twice.} The validation partition selected the
        epoch and fitted the calibration; the test partition, read once, shows
        that this introduced no appreciable bias (Table~\ref{tab:res-test}).
\end{itemize}
